% Page 057

\seite{57}

\abschnitt{Schuljahr 1921–22}
\pstart
Entlassen wurden in diesem Jahre 8 Kinder, welche sich sämtlich für\ob
die Landwirthschaft entschlossen, nämlich Schüler\unsicher: 1\ob
Schüler auf Ostern\unsicher, 1 Knabe und 3 Mädchen; 15 Mädchen\unsicher

\pend
\abschnitt{Lehrerwechsel}
\pstart

Am 1.\ April wurde der bisherige Lehrer Platz in\ob
hiesiger Eigenschaft in die umfassendere Schule nach\ob
Karolich\unsicher{} (bei Daun) versetzt.  Für die hiesige\ob
Schule wurde der Schulamtsbewerber Veit ernannt.

Nachdem die schwierigen Besoldungs- und Gehaltsangelegenheiten gegen\ohb
wärtiger Verhältnisse die hiesige Stellung, infolge der neuen\ob
15.\ Mai 1921  Bestimmungen verlassen mußte, wurde der Schulamts\ohb
bewerber Veit, nach der Verordnung vom 1.\ April\ob
bis 1.\ Mai dienstfrei gelassen. Laut Regierungsbeschluß vom\ob
7.\ April 1921 wurde Lehrer Kneupp, seit 28.\ 5. 19\unsicher{} an\ob
der Schule in Fließem\unsicher, auf den 1.\ Mai als nunmehriger\ob
aufgestellter Lehrer an die hiesige Stelle berufen.\ob
Möge sein Wirken\unsicher{} für Schule und Ge\ohb
meinde gedeihlich sein!

Am 13.\ Mai fand die Firmung in feierlicher Stimmung statt.\ob
Unsre Gemeinde umfaßte 15 Jünglinge,\ob
11 Mädchen, 4 Knaben\unsicher, von denen etwa 2 aus der\ob
Schule schon entlassen waren. Festlich geschmückt\ob
und feierlich bewegte sich der Zug nach althergebrachter Weise\unsicher\ob
in festem Gange zum Orte, über die Geistlichen\ob
Reden nebstdem bei der Andacht die betreffenden\ob
Firmlinge und Geistlichen\unsicher.

24.\ Mai  Vom 14.\ bis einschließlich 23.\ Mai dauerten\ob
die achttägigen Pfingstferien, die von günstig\ob
schönem Wetter begleitet waren.
\pend
